\hspace{3mm}

\centerline{\bf \Huge Делимость+остатки}

\centerline{\bf \large Базовый уровень}

\begin{flushright}
        \scriptsize{}
\end{flushright}

\textbf{Очень важный факт!} \textit{Если каждое из двух целых чисел делится на число $d$, то их сумма, разность и произведение тоже делится на $d$.}

\textbf{Следствие.} \textit{Сумма, разность и произведение любого количества целых чисел делящихся на целое число $d$ тоже делится на $d$.}

\problem{0.1}
Докажите очень важный факт и его следствие.

\problem{0.2} Как вы думаете, среди четырёх последовательных натуральных чисел будет ли хотя бы одно делиться  а) на 2?  б) на 3?  в) на 4?  г) на 5?


\problem{1}
Мафиози Феонардо зашел в свой любимый ресторан Друзья и заказал на всю группу три большие пиццы, шесть картошек фри и двеннадцать бутылок кока колы. "C вас 2150 рублей,"\ --- сказал официант, после чего Феонардо молча достал револьвер. Как он догадался, что его пытаются обсчитать? 

\problem{2}
Миша расставил в клетках прямоугольника $8\times13$ натуральные числа так, что сумма чисел в каждой восьми строк равно 1000. Может ли сумма в каждом из 13 столбцов делиться на 13?

\problem{3}
Докажите, что произведение любых пяти последовательных чисел делится а) на 30; б) на 120.

\problem{4}
У Аюша много-много карточек, на каждой из которых написано число $11, 33, 132$ или $1001$. Он составил из таких карточек несколько натуральных чисел (например, из карточек $11$ и $1001$ можно составить число $111001$) и сложил их всех. Могла сумма быть равно $10^{2024}$?

\problem{5}
В столовой 4 школы продают вкусные булочки. Прозвенел звонок и орда голодных школьников ринулась в столовую. Сперва в столовую прибежали несколько мальчиков и выстроились в очередь. Затем между каждыми двумя мальчиками в очередь влезло по одной девочке. Потом между каждыми двумя детьм влезло еще по одной девочке. Потом между каждыми двумя соседними девочками влезло еще по одному мальчику. Могло ли оказаться, что теперь в очереди стоит 2021 человек?


\newpage

\hspace{3mm}

\centerline{\bf \Huge Делимость+остатки}

\centerline{\bf \large Продвинутый уровень}

\begin{flushright}
        \scriptsize{}
\end{flushright}

\textbf{Определение.} Целые числа \( a \) и \( b \) сравнимы по модулю \( m \), если
\begin{itemize}
    \item они имеют одинаковые остатки при делении на \( m \);
    \item их разность делится на \( m \) (что равносильно предыдущему условию).
\end{itemize}

\textbf{Свойства сравнений.}
\begin{enumerate}
    \item Если \( a \equiv b \pmod{m} \) и \( c \equiv d \pmod{m} \), то \( a \pm c \equiv b \pm d \pmod{m} \).
    \item Если \( a \equiv b \pmod{m} \) и \( c \equiv d \pmod{m} \), то \( ac \equiv bd \pmod{m} \).
    \item Если \( a \equiv b \pmod{m} \), то \( a^n \equiv b^n \pmod{m} \).
    \item Если \( ka \equiv kb \pmod{m} \) и \( (k, m) = 1\), то \( a \equiv b \pmod{m} \).
\end{enumerate}

\problem{1} Найдите остаток от деления

а) \( 47^{101} \) на 46 и 48;

б) \( 2019 \cdot 2020 \cdot 2021 \cdot 2022 \) на 2023;

в) \( 47^{101} \) на 31;

г) \( 9^{2020} + 13^{2020} \) на 11;

д) \( 51! + \frac{102!}{51!} \) на 103.


\problem{2} Докажите, что \( 1^{101} + 2^{101} + 3^{101} + \ldots + 2023^{101} \) делится на 2024.


\problem{3} У числа \( 2024^{2024} \) нашли сумму цифр. У результата опять нашли сумму цифр, и т.д., пока не получилась цифра. Что это за цифра?

\problem{4} Натуральные числа \( a \) и \( b \) таковы, что \( N = (84a + 17b)(17a + 84b) \div 101 \). Докажите, что \( N \ \vdots \ 10201 \).

\problem{5} 
Существует ли степень двойки, из которой перестановкой цифр можно получить другую степень двойки?